\section{Preparing exact data and a double-null template}
\label{pk:section}

Fix a smooth datum $d\in\Ueps$, a branch tube together with all its
witnesses from \Cref{lem:bg-branch-tube}, and a finite integer $m\ge10$.
After interchanging the optical labels if necessary, its central
track has $u\to u_*\in\R$ and $v\to\infty$. All constants in this
section may depend on this fixed datum and tube. Write $\mathcal N$ for its seed
neighborhood. Use the witnessing constants $u_2>u_*$, $c_2>0$, and
$c_0>0$, interval $J$, observation boxes $\mathcal Q_T$, launch patch,
and unbounded sequence of $T$ throughout. The finite experiment uses
\begin{equation}\label{pk:domain}
 \mathcal D_T=\left\{0\le t\le L_T:=\frac{T+u_*}{2}+c_0,
 \quad u+\frac{c_0t}{T}\le u_2,
 \quad v+\frac{c_0t}{T}\le T+c_2\right\}.
\end{equation}
We construct on small enlargements and restrict inside their margins.
The tube, the observation interval $v\in T+J$, and its angular
footprint fit these margins along an unbounded sequence of $T$.
The slice $x$-intervals have a positive lower length bound. Put
\begin{equation}\label{pk:scales}
 \alpha=\kappa/4,\qquad \amp=e^{-\alpha T},
 \qquad \lambda=e^{\zeta T}.
\end{equation}
The small constant $\zeta>0$ is chosen only after all finite base
orders and rate tolerances. In particular, $\amp$ is substantially
larger than the square root of the lapse in the observation region.

Use the observation boxes $\mathcal Q_T$ and their fixed enlargements
from \Cref{lem:bg-branch-tube}. They contain all hitting positions of
the seed-neighborhood closure $\overline{\mathcal N}$ at $v\in T+J$. Their feet lie in
one compact subpatch of the admissible launch patch, with a common
interior margin along the selected sequence. Choose the launch bump
below to be at least one on a neighborhood of this entire foot set.

\begin{proposition}[Packet preparation]\label{prop:packet-preparation}
For every finite $m\ge10$, finite jet budget $N$, positive loss
tolerance $\eta_{\rm out}$, and $0<\nu<1$, one can choose $\zeta>0$ and
finite construction orders so that the following objects exist
for arbitrarily large $T$ along the indicated sequence.
\begin{enumerate}
\item There are exact smooth complete vacuum data $d_T\in\Ddata$
on the whole bridge with $p_m(d_T-d)\to0$. They agree with $d$ on
the other end, the connecting region, and an expanding protected
ball on the selected end. The permitted exterior correction tail
is retained.
\item An exact exterior evolution of $d_T$ reaches the required
part of $\mathcal S$, with overlapping collars and strict outflow
margins. In the fixed background wave-map description its entry
spacetime jets differ from $g+h^p$ by $O(\amp^2e^{\eta_{\rm out} T})$.
Here
\begin{equation}\label{pk:packet}
 h^p=\amp\,\operatorname{tr.rev.}_g\operatorname{Re}
       \left(e^{i\lambda v}\sum_{j=0}^{J_*}\lambda^{-j}p_j\right).
\end{equation}
The leading coefficient is real, trace-free, and annihilates
$L_0$. On the entire observation box $\mathcal Q_T$,
\begin{equation}\label{pk:screen-lower}
 |(p_0)_{AB}|_\gamma\ge e^{-o(T)}.
\end{equation}
At every fixed required base order the ordinary mixed coordinate jets
of the normalized amplitude components are $e^{o(T)}$, and the linearized Ricci and gauge errors have
arbitrarily high inverse-frequency accuracy.
\item There is a metric $\bar g$ exactly of double-null form, with
fields $\bar\gamma,\log\bar{\mathcal A},\bar b$ and their exact
geometric derived fields. Through order $N$, its differences from
the background fields, including $q^{-1}\partial_vb$, are
$O(\amp e^{\eta_{\rm out} T})$. Its defects in the null structure and
angular equations \eqref{dn:D2}--\eqref{dn:D3} are
$O(\amp^2e^{\eta_{\rm out} T})$. If the actual exterior induced data are
written in double-null coordinates initialized by the background
values on $\mathcal S$, their full entry jets differ from this
template by the same quadratic bound.
\end{enumerate}
All entry discrepancies and all template changes and defects vanish
for $v<(1-\nu)T$. Any additional finite initial or residual jet
buffers may be included in $N$. Positive frequency powers at these
base orders are independent of the later orders used to improve
linear accuracy. At the finitely many high entry orders ultimately
requested by \Cref{thm:dn-comparison}, the actual fields have
$e^{o(T)}$ bounds.
\end{proposition}

For the proof, choose an internal loss
$0<\eta<\min\{\eta_{\rm out},\alpha/4\}$. The bootstrap estimates
below are obtained with this smaller loss, after splitting it among
the finitely many steps, and then imply the stated bounds with
$\eta_{\rm out}$. This distinction permits any prescribed positive
output tolerance.

\subsection{The homogeneous linear hierarchy and its footprint}

Complex Gaussian beams and their superpositions were developed by
Ralston, Tanushev, and Liu--Runborg--Tanushev
\cite{Ralston1982,Tanushev2008,LiuRunborgTanushev2013}.
The companion construction used here also enforces the Einstein
wave-gauge constraints and supplies the fixed-order bounds on
intervals growing with $T$. We retain the hierarchy and matching
estimates needed for the larger amplitude.

We use the amplitude-independent constrained construction of
\cite[\CompanionRef{signed}{bg:packet}, \CompanionRef{signed}{bg:amplitude-independent}]{KerrCOne}, whose hypotheses are
the background coefficient estimates and a branch tube; it expressly
requires no bound on curvature in an observer's parallel frame.
Its proof uses the hierarchy and beam interfaces of
\cite[\CompanionRef{foundation}{pk2:hierarchy}, \CompanionRef{foundation}{pk2:escape}, \CompanionRef{foundation}{pk2:beam-proposition}]{KerrCTwo}.
We specify the part needed for a signal on a whole footprint.

Let $P$ be the background linearized reduced vacuum operator on
trace-reversed two-tensors, $C_g=\divg_g$, and $P_1$ the induced
one-form wave operator. The Ricci-flat Bianchi identity is
$C_gP=P_1C_g$. For
\[
 k_v=g^{-1}\dd v=-a^{-1}L_0,\qquad
 K_vp=p(k_v,\cdot),\qquad
 \mathcal T=2\nabla_{k_v}+\square_gv,
\]
the coefficients solve
\begin{equation}\label{pk:hierarchy}
 i\mathcal T p_j+Pp_{j-1}=0,\qquad
 iK_vp_j+C_gp_{j-1}=0,\qquad p_{-1}=0.
\end{equation}
Work on $-v\le u\le u_f$, where $u_f>u_2$. Multiplication of the
homogeneous transport by $-a/2$ gives
$\nabla_{L_0}+\tfrac12\tr_\gamma\chi_0$. In the normalized slots
$(a^{-1/2}L_0,a^{-1/2}L_1,\partial_\theta)$ its constant diagonal
weights are $\kappa(n_0-n_1)/2$: here $n_i$ counts the arguments of
type $a^{-1/2}L_i$ in a covariant two-tensor component, so
$n_0+n_1\le2$. The normalized connection formulas
in \cite[\CompanionRef{foundation}{pk2:coefficients}]{KerrCTwo} make the hierarchy triangular from
higher weights to lower weights. Nonnegative weights are prescribed
at $u=-v$, with the components containing an $L_0$ argument fixed
by the gauge equation;
negative weights $1A,11$ are prescribed at $u_f$ and transported
backwards. The sign of each weight is favorable in its prescribed
direction. Higher differentiated coefficients multiply already
controlled lower jets, so each finite transport induction costs
only $e^{o(T)}$. The angular characteristic and inverse
characteristic jets have this same bound, including derivatives
of the moving initial edge $u=-v$.

For the free leading angular data, orthonormalize a fixed smooth local
coframe using the entry screen metric, obtaining $e^1,e^2$. Uniform
low nondegeneracy and the fixed higher coefficient bounds give
$e^{o(T)}$ jets for this coframe at every fixed order. Use the real
trace-free unit polarization
$2^{-1/2}(e^1\otimes e^1-e^2\otimes e^2)$, multiplied by the bump
that is at least one on a neighborhood of the complete foot set.
Take zero free angular data at higher levels.
The initial polarization exists on the one small chart patch; no
global choice of a nonvanishing trace-free tensor on $\Sph$ is
required. The leading transport is real. It preserves trace and
annihilation of $k_v$, and acts on the positive screen quotient by
isometric parallel transport times a real nonzero scalar. The
scalar and its reciprocal are $e^{o(T)}$, since its logarithm is
an integral of the expansion coefficient. Thus the positive bump
gives \eqref{pk:screen-lower} at every point whose foot lies in
the chosen foot set, uniformly throughout $\mathcal Q_T$.
There is no phase adjustment at a single observation point.

Gauge compatibility follows directly by conjugating
$C_gP=P_1C_g$ with $e^{i\lambda v}$: once preceding gauge defects
vanish, the next defect satisfies homogeneous one-form transport
with zero initial value. Terminal prescriptions for negative
weights create no additional condition, since $K_v$ contracts
precisely an $L_0$ argument. Smooth zero extension away from the launch
band is therefore allowed. For any fixed base order, truncating
\eqref{pk:hierarchy} far enough gives scaled reduced-wave and gauge
residuals $\amp\lambda^{-M}e^{o(T)}$ for any prescribed $M$.

The exterior transfer of \cite[\CompanionRef{signed}{bg:packet}, \CompanionRef{signed}{bg:amplitude-independent}]{KerrCOne} supplies a real
lab-coordinate tensor $h^B$, including the factor $\amp$, with
\begin{equation}\label{pk:beam-base}
 \|h^B\|_{C^s}\le\amp\lambda^{s+3/2}e^{o(T)}.
\end{equation}
Its wave, gauge, and entry-matching residuals can simultaneously be
$\amp\lambda^{-M}e^{o(T)}$. The geometric hypotheses of that lemma
hold here: the prepared optical covectors have escaping backward
rays; the entry, horizon, and initial-slice crossings are
transverse; the fixed-depth collar has uniform low nondegeneracy;
and all fixed higher coefficient and flow jets are subexponential.
These are coefficient properties of the fixed background from
\Cref{prop:bg-package}. The fixed small launch patch is chosen
using only these low bounds. It need not shrink as approximation
orders increase. For each prescribed relative annulus tolerance,
a larger fixed far threshold and then large $T$ place the initial
support in that thin annulus, uniformly for all centers in the
patch; this is the backward-ray estimate in
\cite[\CompanionRef{foundation}{pk2:escape}]{KerrCTwo}.

We recall why the positive exponent in \eqref{pk:beam-base} is
independent of later accuracy. The constrained beams are integrated
over a fixed three-dimensional center patch with normalization
$(\lambda/2\pi)^{3/2}$. Each of $s$ differentiations costs at most
$\lambda$, while every finite phase and amplitude coefficient is
$e^{o(T)}$ after its displayed inverse powers have been extracted.
Taking absolute center integrals gives \eqref{pk:beam-base}.
For each beam center $c$, let $P_c,Q_c\in\mathbb C^{3\times3}$ be
the position and momentum differential blocks of the Hamiltonian ray
map acting on the initial complex phase graph. If $M_0$ is its
symmetric initial phase Hessian with positive imaginary part, then
$P_c(0)=\Id$ and $Q_c(0)=M_0$. The imaginary Hessian remains positive
through real caustics by the symplectic identity
$P_c^*Q_c-Q_c^*P_c=2i\operatorname{Im}M_0$; it gives a
subexponential inverse of the complex position projection.
At a fixed transverse Taylor degree $d$, homogeneous amplitude-jet
transport in ray-comoving coordinates is the central amplitude
transport. Returning to fixed coordinates also applies the $d$-fold
symmetric tensor power of $J_c(t,s)^{-{\mathsf t}}$, where
$J_c(t,s)=P_c(t)P_c(s)^{-1}$. The bounds for $P_c$ and $P_c^{-1}$
keep this factor subexponential for every fixed $d$. Thus same-degree
jet couplings remain controlled when the finite accuracy order is raised.
An uncancelled Taylor degree gains
$\lambda^{-1/2}e^{o(T)}$ by Gaussian integration. Increasing the
finite Taylor and transport orders therefore improves only
inverse accuracy. The entry Gaussian expansion is triangular in
amplitude level, and its new free contraction-kernel value matches
the corresponding $p_j$. These are precisely the fixed-base
assertions proved in \cite[\CompanionRef{foundation}{pk2:beam-proposition}]{KerrCTwo}. Linearity makes
them valid at the amplitude \eqref{pk:scales}.

\subsection{Exact constraints and exterior propagation at the new amplitude}

Localized deformation of the vacuum constraints was developed by
Corvino--Schoen and Chru\'sciel--Delay
\cite{CorvinoSchoen2006,ChruscielDelay2003}. Here the correction fixes
an inner ball and retains a noncompact decaying tail. Its support and
weighted bounds must remain explicit, since they determine both
closeness of the complete data and the first time the perturbation can
reach the interior.

In the selected exterior, $t_e$ denotes exterior time, $z$ lab
position, and $C(t_e)$ the drift center; its co-moving radius is
$r_{\rm co}=|z-C(t_e)|$.
Write $t_T\asymp T$ for the exterior time corresponding to the
launch cylinder and $z_*=C(t_T)$ for its lab drift center. Choose
a small relative gap $\beta>0$ with
$s_*:=\sup|C'|<1-2\beta$, and set
\begin{equation}\label{pk:protected}
 \rho_T=(1-2\beta)t_T,\qquad |z-z_*|<\rho_T.
\end{equation}
The thin launch annulus lies outside this protected ball with a
margin. In the variables $y=(z-z_*)/\rho_T$, keep Cartesian metric
components unchanged and multiply second-fundamental-form
components by $\rho_T$. Both constraints are then multiplied by
$\rho_T^2$. Extending the unperturbed rescaled datum smoothly
inside the unit ball, without imposing a constraint there, gives
a perturbation of flat data of size $O_s(T^{-1})$ at every fixed
symbol order.

Here is the exact constraint-solver interface being used. For
$s\ge3$, let
\[
 \|f\|_{H^j_\ell}=\|f\|_{H^j(B_2)}+
  \sup_{R\ge1}R^\ell
       \|f(R\,\cdot)\|_{H^j(\{1<|y|<2\})},\qquad
 X^s=H^{s+2}_1\oplus H^{s+1}_2,\quad Y^s=H^s_4.
\]
Here $H^j_\ell$ is the class of locally $H^j$ fields with finite
displayed norm. In particular, the weight-one and weight-two factors
include fields with saturated $|y|^{-1}$ and $|y|^{-2}$ decay at
infinity; these classes are larger than the closures of compactly
supported fields.
\cite[\CompanionRef{foundation}{pk:constraint-nonlinear}, \CompanionRef{foundation}{pk:constraint-linear}]{KerrCTwo} give a bounded
operator $\mathcal K:Y^s\to X^s$ preserving exterior support, with
$P_{\rm flat}\mathcal Kz=z$ for exterior-supported $z$, where
$P_{\rm flat}$ is the flat constraint linearization. They also give a contraction for the nonlinear
constraints when one finite background-plus-defect base norm is
small. For the smooth compact beam deformation and the smooth
background with all-order symbol bounds used here, the output is
smooth, zero inside the protected ball, and has all-order symbol
tails of weights one and two. The iteration sources may themselves
have noncompact exterior-supported tails. There is no
moment condition on the original defect and no fixed-charge
requirement. Smoothness at higher orders requires bounded higher
symbol norms, not their smallness.

Take $s$ large enough for $m$ and all requested comparison buffers.
Let $d_B$ be the rescaled beam seed, and let $f_B$ be its exterior
constraint defect extended by zero inside the ball. Thus $f_B=\Phi(d_B)$
outside the ball; the artificial interior need not satisfy the constraints.
Linearize the constraint map at the actual rescaled vacuum background
on the exterior. Gauss--Codazzi expresses this first variation as a
linear combination of the reduced-wave residual and one derivative
of the linear gauge defect, with smooth background coefficients.
Include gauge-defect jets through $s+1$ among the fixed base orders
before choosing $\zeta$. Taylor expansion around that background then gives
\begin{equation}\label{pk:exterior-defect}
 \|f_B\|_{Y^s}
 \le \operatorname{poly}(T)e^{o(T)}
      \bigl(\amp\lambda^{-M}+\amp^2\lambda^{P_*}\bigr).
\end{equation}
Here and below $P_*$ depends on the fixed base orders alone.
The flat operator is reserved for the correction iteration below;
the linear cancellation in \eqref{pk:exterior-defect} uses the actual
background.
The rescaled seed base norm is
$O_s(T^{-1})+\operatorname{poly}(T)\amp\lambda^{P_*}e^{o(T)}$.
Choosing $\zeta P_*<\alpha/2$ makes it small. The inverse-metric
and contraction balls of the cited proposition consequently
apply for large $T$. Equivalently, solve
\[
 z=-f_B-N_{d_B}(\mathcal Kz),\qquad
 N_{d_B}(w)=\Phi(d_B+w)-\Phi(d_B)-P_{\rm flat}w.
\]
The Lipschitz constant is at most $1/2$, and the correction has
the size of \eqref{pk:exterior-defect}. The defect is smoothly zero
near the inner boundary, and the support-preserving inverse and
nonlinear iteration keep the correction zero on the entire inner
side. Since $P_{\rm flat}\mathcal Kz=z$, the corrected constraint
is $\Phi(d_B+\mathcal Kz)=\Phi(d_B)-f_B$, hence zero outside the ball.
On the inner side $N_{d_B}(\mathcal Kz)=0$ because the correction and
all its jets vanish. Undoing the rescaling and restoring the original
vacuum datum inside therefore gives smooth exact data on the complete
bridge.

The metric stays positive and uniformly comparable to the original
complete metric. Weighted differentiation after translation and
rescaling costs only powers of $T$. For example, the seed contributes
at most $\operatorname{poly}(T)\amp\lambda^{m+5/2}e^{o(T)}$ to
$p_m$, and the correction contributes the same polynomial type of
loss times \eqref{pk:exterior-defect}. Thus $p_m(d_T-d)\to0$.
Every higher weighted seminorm of each individual $d_T$ is finite
by the all-order symbol assertion. The exterior tail and any charge
change remain part of $d_T$.

Initialize the exact exterior wave-map evolution with target $g$.
Choose initial lapse and shift as for $g+h^B$; the corrected second
fundamental form and exact gauge determine the remaining normal
metric jets by a uniformly noncharacteristic algebraic system.
Its linearization matches the constrained beam, so its discrepancy
has size \eqref{pk:exterior-defect}. Further normal jets follow
from the reduced equations. Agreement on the protected region is
exact.

We now justify the new amplitude in the exterior propagation.
For a lab-component metric change $H$, the reduced equation has
the form
\begin{equation}\label{pk:exterior-operator}
 \mathcal R_e(H)_n=(\mathcal L_eH)_n+
 \bigl((g+H)^{-1}-g^{-1}\bigr)^{\rho\sigma}
       \partial_\rho\partial_\sigma H_n+
 \mathcal Q_n(H,\partial H),
\end{equation}
where $\mathcal Q$ starts quadratically and has total derivative
order at most two, with fixed differentiated coefficients
$e^{o(T)}$. On the uniformly nondegenerate exterior domain,
$\mathcal R_e(h^B)$ has the size in
\eqref{pk:exterior-defect}; there is no small interior lapse in
this assertion. Write $H=h^B+F$ and bootstrap sufficiently high
mixed lab-coordinate error energies by $e^{-\alpha T}$.
Retain the actual metric $g+h^B+F$ in the highest-order wave
operator. Every other nonlinear error term has an additional
packet or error factor. Put its highest error derivative in
$L^2$ and its lower factors in $L^\infty$ using a fixed Sobolev
buffer. These factors are exponentially small, so the nonlinear
energy rate is exponentially small.

The individual exterior estimate proved in
\cite[\CompanionRef{signed}{prep:exterior-rate-section}, \CompanionRef{signed}{prep:amplitude-independent-exterior}]{KerrCOne}
uses finite shrinking computing slices
\begin{equation}\label{pk:outer-cutoff}
 r_{\rm co}<R_{\rm out}(t_e)
 =C_{\rm out}(1+T)+V_{\rm out}(t_T+C_2-t_e),\qquad
 V_{\rm out}>\sup_{\rm causal}|\dd r_{\rm co}/\dd t_e|.
\end{equation}
Their radii are $O(T)$ on the finite comparison interval, and the side
is strictly outgoing. The physical correction tail outside this
computing region remains part of the complete initial datum. On these
slices the estimate for square-root energies at exact orders $j$ reads
\begin{equation}\label{pk:exterior-energy}
 E_j'\le r_e E_j+C_je^{o(T)}\sum_{k<j}E_k+\mathcal F_j,
 \qquad \mathcal F_j\le
 \operatorname{poly}(T)e^{o(T)}
 (\amp\lambda^{-M}+\amp^2\lambda^{P_*}).
\end{equation}
For precision, its arbitrarily small integrated rate is available
at this amplitude for the following reason. The actual scalar
stress energy, with future slice normal and a small positive mass,
has nonnegative outflow flux. Its same-order linear rate uses only
first and second lab metric derivatives, which are small at large
co-radius and uniformly bounded elsewhere. Higher background
derivatives multiply lower exact-order energies. The first-arrival
estimate, valid also for the noncompact correction since it is zero
in the protected ball, is
\begin{equation}\label{pk:first-arrival}
 (1-2\beta)t_T\le (1+\delta_R)t_e+O(R)
                              +s_*|t_T-t_e|,
 \qquad \delta_R\longrightarrow0.
\end{equation}
Indeed before its first hit of a fixed large co-radius $R$, a
causal path from the changed support has lab speed at most
$1+\delta_R$; its endpoint is within $O(R)$ of $C(t_e)$.
The fixed gap $1-s_*>0$ in \eqref{pk:first-arrival} implies that
the change cannot see the bounded-rate region before
$t_e\ge t_T-\beta'T$, where $\beta'>0$ can be made arbitrarily
small by decreasing $\beta$ and increasing $R$. The localized
beam obeys the same conclusion. Thus
\[
 \int_0^{t_T+C_2}r_e\,\dd t_e\le\eta' T
\]
for any prescribed $\eta'>0$. The bounded compact rate is
independent of the chosen far threshold. Increasing that threshold
later for the thin annulus therefore causes no conflict. The
choices also protect the entire earlier entry interval
$v<(1-\nu)T$.

Remove the common rate from \eqref{pk:exterior-energy}.
Iteration of the finite lower-triangular hierarchy introduces only
a polynomial of fixed degree in $Te^{o(T)}$, hence another
$e^{o(T)}$ factor. Allocating the finitely many rate losses in
advance, then choosing $\zeta$ small and $M$ large, gives
\begin{equation}\label{pk:wave-entry}
 \|g_T^{\rm ext}-g-h^p\|_{C^N(\mathcal S)}
                         \le\amp^2e^{\eta T}.
\end{equation}
The bound first holds with $h^B$ on the exterior; the fixed-depth
coordinate transition, spacetime trace with finite buffers, and
linear entry matching replace it by $h^p$. It strictly improves
the bootstrap since $2\alpha-\eta>\alpha$ after choosing losses
small. Local reduced-wave restart inside enlarged outflow margins
gives the entire finite exterior. Constraint and gauge propagation
give exact vacuum. These continuation statements are the ones
used in \cite[\CompanionRef{foundation}{in:continuation}, \CompanionRef{foundation}{cmp:development}, \CompanionRef{foundation}{pk2:exterior-proposition}]{KerrCTwo}.
This argument is the individual exterior comparison; it uses no
deep wave-coordinate nonlinear estimate.

\subsection{Linear optical coordinates and the relative shift}

The exact evolution from the entry cylinder will use two optical
functions and angles having the \emph{same values} as the background
ones on $\mathcal S$. Choose the neighboring eikonal roots and
transport the angles so that the shift in direction $1$ is zero.
The entry surface consequently retains the labels $u+v=0$.
Only the linearized coordinate change is needed to construct the
comparison template.

\begin{lemma}[Linear optical gauge]\label{lem:packet-linear-gauge}
Let $Z=(z^0,z^1,z^\theta)$ be the variations of the coordinate
functions, with $Z=0$ on $\mathcal S$, for the input $h^p$.
In unnormalized components $h^p_{ij}=h^p(L_i,L_j)$ and
$h^p_{iA}=h^p(L_i,\partial_A)$, they obey
\begin{equation}\label{pk:linear-optical}
 \begin{aligned}
 \partial_vz^0&=-\frac{h^p_{11}}{2a},&
 L_0z^1&=-\frac{h^p_{00}}{2a},\\
 \partial_vz^\theta
 &=b\partial_vz^0+a\gamma^{-1}\dd_\theta z^0
                         +\gamma^{-1}h^p_{1\cdot}.
 \end{aligned}
\end{equation}
For every fixed number of angular derivatives and $r$ total
longitudinal derivatives,
\begin{equation}\label{pk:coordinate-gain}
 |\partial_{u,v}^{r}\partial_\theta^sZ|
                      \le\amp\lambda^{r-1}e^{o(T)}.
\end{equation}
The fixed-new-label metric variation is $k=h^p-\Lie_Zg$.
Its field variations $k_\gamma,k_H,k_b$, where $H=\log\mathcal A$,
and the relative derivative $q^{-1}\partial_vk_b$ satisfy
\begin{equation}\label{pk:relative-linear}
 \|(k_\gamma,k_H,k_b,q^{-1}\partial_vk_b)\|_{C^s}
                        \le\amp\lambda^{P_s}e^{o(T)}
\end{equation}
at every fixed required order. The exponents $P_s$ depend only on
that order, not on subsequent inverse-accuracy choices.
\end{lemma}

\begin{proof}
The inverse metric is
\[
 g^{-1}=-a^{-1}(L_0\otimes L_1+L_1\otimes L_0)
                        +\gamma^{AB}\partial_A\otimes\partial_B.
\]
Linearization of $g^{-1}(\dd u,\dd u)=0$ gives
$-a^{-2}h^p_{11}-2a^{-1}\partial_vz^0=0$; the equation for $v$
is analogous. For the angle condition
$g^{-1}(\dd u,\dd\theta^A)=0$, the metric variation contributes
$a^{-1}\gamma^{AB}h^p_{1B}-a^{-2}b^Ah^p_{11}$.
The coordinate variations contribute
$\gamma^{AB}\partial_Bz^0-a^{-1}b^A\partial_vz^0
-a^{-1}\partial_vz^A$.
Substitution of the first eikonal equation gives the last equation
in \eqref{pk:linear-optical}, including its sign and factor.

To see the gain in \eqref{pk:coordinate-gain}, a typical term in
the first equation has the form
$\amp e^{i\lambda v}f(u,v,\theta)$ with fixed coefficient jets
$e^{o(T)}$. At fixed $u,\theta$,
\begin{equation}\label{pk:oscillatory-integral}
 \begin{split}
 \int_{-u}^{v}e^{i\lambda s}f(u,s,\theta)\,\dd s
 ={}&\frac{e^{i\lambda v}f(u,v,\theta)
          -e^{-i\lambda u}f(u,-u,\theta)}{i\lambda}\\
 &-\frac1{i\lambda}
           \int_{-u}^{v}e^{i\lambda s}\partial_sf(u,s,\theta)\,\dd s.
 \end{split}
\end{equation}
The interval has length $O(T)$; its polynomial length is absorbed
in $e^{o(T)}$. Differentiation in $u$ differentiates the lower
endpoint and its coefficient as well as the integral. Each
longitudinal derivative costs at most one additional $\lambda$,
including the phase $e^{-i\lambda u}$ at that endpoint. Angular
derivatives do not differentiate the phase. This proves the
assertion for $z^0$ with arbitrary fixed differentiated orders.
The relative coefficient $h^p_{11}/a$ is a normalized tensor
entry, so no inverse lapse was lost.

For $z^1$, the leading $00$ coefficient is identically zero:
$p_0$ annihilates $L_0$, and trace reversal leaves the same-null
$00$ component unchanged because $g_{00}=0$. Thus $h^p_{00}/a$ already has a factor
$\lambda^{-1}$. Integration along $L_0$ keeps $v$ fixed, with
angular characteristic jets and inverse jets $e^{o(T)}$.
Differentiating its moving initial point $u=-v$ has the same
endpoint rule as above. Its interval length is again $O(T)$.
This proves the bound for $z^1$ without requiring oscillation in
the direction of that integration.

For the angular equation, $h^p_{1A}$ is $\sqrt a$ times a
normalized oscillatory coefficient, so \eqref{pk:oscillatory-integral}
applies. The term $a\gamma^{-1}\dd_\theta z^0$ already has the
inverse-frequency gain. Integrating its length costs only
$e^{o(T)}$, also at differentiated orders. Finally,
\[
 \int_{-u}^{v}b\partial_vz^0\,\dd v
   =bz^0\big|_{-u}^{v}-\int_{-u}^{v}(\partial_vb)z^0\,\dd v,
\]
where the lower value of $z^0$ is zero. All derivatives of
$\partial_vb$ retain their factor $a$, and the remaining
coefficients have subexponential fixed jets. This proves
\eqref{pk:coordinate-gain} for the angles as well.

For explicit field formulas set $W=z^\theta-z^0b$, an angular
vector field. A direct Lie derivative gives
\begin{equation}\label{pk:linear-fields}
 \begin{aligned}
 k_\gamma&=h^p_{AB}-2z^0\chi_0-2z^1\chi_1-\Lie_W\gamma,\\
 k_H&=-a^{-1}h^p_{01}-Z\log a-L_0z^0-\partial_vz^1,\\
 k_b&=-\gamma^{-1}h^p_{0\cdot}
       +\partial_uW+[b,W]-z^1\partial_vb
                                      -a\gamma^{-1}\dd_\theta z^1.
 \end{aligned}
\end{equation}
Indeed $(\Lie_Zg)_{AB}=2z^0\chi_0+2z^1\chi_1+\Lie_W\gamma$,
whereas
$(\Lie_Zg)(L_0,L_1)=-Za-a(L_0z^0+\partial_vz^1)$.
The $0A$ component is
$\gamma_{AB}(\partial_uW+[b,W]-z^1\partial_vb)^B
-a\partial_Az^1$. The linear double-null component conditions
$k_{11}=k_{00}=k_{1A}=0$ follow from
\eqref{pk:linear-optical}. Hence
$k_H=-k(L_0,L_1)/a$ and $k_b=-\gamma^{-1}k_{0\cdot}$, which
give \eqref{pk:linear-fields}. All displayed divisions by $a$
are relative factors: $h^p_{01}/a$ is normalized and
$\partial a/a=\partial\log a$ has the given coefficient bounds.
Thus these three field variations have the claimed fixed
positive frequency powers. This argument alone does not estimate
$q^{-1}\partial_vk_b$; that relative derivative is obtained next.

Let $X_i$ denote the exact linear variation of $\xi_i$ under $k$.
The torsion sum identity gives $X_0+X_1=\dd_\theta k_H$.
The linearized vacuum identities are available with arbitrary
linear accuracy: diffeomorphism covariance on the vacuum
background gives
\[
 D\Ric_g(\Lie_Zg)=\Lie_Z\Ric(g)=0.
\]
The algebraic and differential steps turning Ricci components into
the null equations can lose fixed powers of $q^{-1}$, but only
on the \emph{linear residual}. On \eqref{pk:domain} these powers
are $e^{C_sT}$ for a fixed $C_s$. After $\zeta$ is fixed, increasing
$M$ in $\amp\lambda^{-M}$ makes this residual smaller than any
prescribed exponential at all finitely many required orders.

Use specifically the shift-free direction-$1$ torsion equation
of \cite[\CompanionRef{foundation}{in:vacuum}]{KerrCTwo}. Its linearization is
\begin{equation}\label{pk:linear-torsion}
 (\partial_v+\tau_1)X_0
   =\delta\beta_1+\delta\tau_1\,\xi_1
                     +\tau_1\dd_\theta k_H+\mathfrak r_{
                     \rm lin}.
\end{equation}
Here
\[
 \delta\chi_1=\tfrac12\partial_vk_\gamma,\qquad
 \delta\tau_1=\tr_\gamma\delta\chi_1
                 -\gamma^{AC}k_{\gamma,CD}\gamma^{DB}\chi_{1,AB},
\]
and
\[
 \delta\beta_1=\divg_\gamma\delta\chi_1
 -\dd_\theta\delta\tau_1+(\delta\divg_\gamma)\chi_1.
\]
The operator variation in the last term uses only $k_\gamma$
and one angular derivative of $k_\gamma$. Thus the entire source
in \eqref{pk:linear-torsion} has already been estimated using
\eqref{pk:linear-fields}; it uses neither $X_i$ nor an unknown
relative shift derivative. Also
\[
 \int_{-u}^{v}|\tau_1(u,s,\theta)|\,\dd s
       \le C\int_{-u}^{v}(q+g_1(s))\,\dd s\le C.
\]
At the moving entry edge $q=q_0$, the initial $X_0$ is a finite-jet
expression with no small denominator. It has size
$\amp\lambda^{P_s}e^{o(T)}$. The integrating-factor estimate
therefore gives this bound for $X_0$, and then for $X_1$.
Commutation with $u$ and angular derivatives leaves the same top
transport coefficient; differentiated coefficients multiply lower
already estimated $X_0$ jets. Entry terms from differentiating
$v=-u$ are controlled by the differentiated entry data, with normal
jets also given by the equation. Derivatives in $v$ are recovered
from \eqref{pk:linear-torsion}. This proves the finite-order
assertion without a shifted transport containing an unknown $k_b$.

Finally linearize the defining shift relation. It gives exactly
\begin{equation}\label{pk:relative-shift}
 \begin{split}
 q^{-1}\partial_vk_b={}&\mathcal A\gamma^{-1}(X_1-X_0)\\
 &+\bigl(k_H\Id-\gamma^{-1}k_\gamma\bigr)
                                      q^{-1}\partial_vb.
 \end{split}
\end{equation}
The background relative derivative and all its fixed jets are
subexponential. This proves \eqref{pk:relative-linear}. In
particular the proof never divides an ordinary error estimate
for $\partial_vk_b$ by a small $q$.
\end{proof}

\subsection{Exact template definitions and quadratic defects}

Define the template metric by the following exact assignments:
\begin{equation}\label{pk:template}
 \begin{gathered}
 \bar\gamma=\gamma+k_\gamma,\qquad
 \bar{\mathcal A}=\mathcal A e^{k_H},\qquad
 \bar a=q\bar{\mathcal A},\qquad \bar b=b+k_b,\\
 \bar g=-2\bar a\,\dd u\,\dd v+
 \bar\gamma_{AB}(\dd\theta^A-\bar b^A\dd u)
                         (\dd\theta^B-\bar b^B\dd u).
 \end{gathered}
\end{equation}
For sufficiently large $T$, $\bar\gamma$ is positive definite
and $\bar{\mathcal A}$ is positive with the background low
nondegeneracy bounds. All barred derived fields are their actual
geometric definitions, rather than independent truncated series:
\begin{equation}\label{pk:template-derived}
 \begin{aligned}
 \bar\chi_i&=\tfrac12(\partial_{y_i}+\Lie_{\bar B_i})\bar\gamma,
 &\bar\ell_i&=(\partial_{y_i}+\bar B_i)\log\bar a,\\
 \bar\tau_i&=\tr_{\bar\gamma}\bar\chi_i,
 &\bar\beta_i&=\divg_{\bar\gamma}\bar\chi_i-\dd_\theta\bar\tau_i,\\
 \bar\xi_0&=\tfrac12\left(\dd_\theta\log\bar{\mathcal A}
       -\bar{\mathcal A}^{-1}\bar\gamma
                                      q^{-1}\partial_v\bar b\right),
 &\bar\xi_1&=\tfrac12\left(\dd_\theta\log\bar{\mathcal A}
       +\bar{\mathcal A}^{-1}\bar\gamma
                                      q^{-1}\partial_v\bar b\right),\\
 \bar c_i&=\curl_{\bar\gamma}\bar\xi_j,
 &\bar Q_i&=(K_{\bar\gamma},\bar c_i),\qquad
 \bar m_i=K_{\bar\gamma}-\divg_{\bar\gamma}\bar\xi_j,\quad j\ne i.
 \end{aligned}
\end{equation}
Here $\bar B_0=\bar b$ and $\bar B_1=0$. In particular the torsion
definitions use the relative shift derivative already controlled
by \Cref{lem:packet-linear-gauge}.

Consider the null identities \eqref{dn:D2} and the angular system
\eqref{dn:D3} as finite differential expressions in
\[
 \gamma,\ \gamma^{-1},\ \mathcal A,\ \mathcal A^{-1},\ b,
 \ H=\log\mathcal A,\ V=q^{-1}\partial_vb.
\]
Their only radial lapse factors are nonnegative powers of $q$.
For example, the cross-shape sources are
$q\mathcal A\nabla\xi$, $q\mathcal A K_\gamma\gamma$,
$q\mathcal A\xi^2$, and products of the two shapes; the lapse
cross equation has $q\mathcal A K_\gamma$,
$q\mathcal A\xi^2$, and shape products. The angular pair has
principal coefficient $q\mathcal A/2$; its other terms are
$q\mathcal A(K_\gamma\xi+\xi\nabla\xi+\xi^3)$ and
$\chi_i\nabla\chi_j+\chi_j\nabla\chi_i$. The remaining angular
equations use $\chi_i(K_\gamma,\nabla\xi)$ and
$\nabla\chi_i\xi$. The shift identity is
$\partial_vb=q\mathcal A\gamma^{-1}(\xi_1-\xi_0)$.
The trace equation can be written without a lapse derivative as
$(L_i+\kappa)(\tau_i/\mathcal A)=-\mathcal A^{-1}|\chi_i|^2$.
Angular curvature and div--curl operations involve only the sphere
metric and its derivatives. Finally, ordinary differentiation of
$q$ preserves its factor. This list accounts for all the factor
types in these equations and their finite differentiated versions.

It follows that substituting \eqref{pk:template}--\eqref{pk:template-derived}
into these expressions is a smooth finite-jet operation on a
uniform low nondegeneracy set. Its background term vanishes and
its linear term is the arbitrarily accurate linearized vacuum
identity proved above. Its remaining terms contain at least two
linear increments, so at each required order the defect is
\begin{equation}\label{pk:template-defect}
 O\left(\amp\lambda^{-M}e^{C_NT+o(T)}
                   +\amp^2\lambda^{P_N}e^{o(T)}\right).
\end{equation}
This separation is essential: fixed inverse-$q$ losses have been
allowed on the first, linear term only. There is no such loss on
the quadratic term, by the displayed relative-field formulas.
Choose $\zeta P_N$ below the allocated loss tolerance, and then
choose $M$ large. Formula \eqref{pk:template-defect} becomes the
quadratic bound in \Cref{prop:packet-preparation}. The same
Taylor argument bounds all template-background field differences
by $\amp e^{\eta T}$.

On the entry cylinder $q=q_0$ and all eikonal roots and angular
transports are uniformly noncharacteristic at low order. Their
initial spacetime jets are smooth finite-jet functions of the
exterior metric jets and of the prescribed coordinate values.
Their linearization is \eqref{pk:linear-optical}, and changing
metric components to the new labels has linearization
$h^p-\Lie_Zg$. Thus \eqref{pk:wave-entry}, with a finite buffer
and initially smaller loss, gives actual-minus-template entry
jets of size $\amp^2e^{\eta T}$. Nonlinear coordinate terms are
quadratic finite-jet terms at this fixed depth. They do not
require a comparison of the two gauges deep in the wedge.
The new coordinates retain the same entry labels. The first-arrival
property, causal monotonicity in double-null coordinates, and
uniqueness of the optical transports show that exact changes and
template changes vanish in $v<(1-\nu)T$ wherever the local exact
double-null evolution exists.

The finite choice order can now be stated without a circular
increase of frequency losses. First choose the target observation
and bootstrap orders. The interpolation and continuation argument
in \Cref{thm:dn-comparison} requests finitely many higher entry and
residual orders; include them and $m$ in the base budget $N$.
Next choose all energy loss tolerances and the support tolerance
$\nu$, including the exterior first-arrival choices. Choose
$\zeta>0$ small for every positive frequency power at these base
orders. Only then enlarge the hierarchy, beam, phase-Taylor, and
matching accuracy orders to absorb every linear inverse-$q$ loss.
At the fixed high entry orders, the packet changes are then
exponentially small, and the background jets are $e^{o(T)}$;
so the actual high entry jets are $e^{o(T)}$ as asserted. Finally
take $T$ large. All choices are finite, and each $d_T$ is smooth
at all remaining orders by the cited constraint and vacuum
regularity results. This proves \Cref{prop:packet-preparation}.
